% Metodologia de validação
\begin{frame}{Metodologia de validação}
  \small
  \begin{columns}[T,onlytextwidth]
    \begin{column}{0.31\textwidth}
      \begin{block}{1. Análise teórica}
        \begin{itemize}
          \item aplicação da LKC;
          \item obtenção de $H(s)$;
          \item substituição $s=j\omega$;
          \item cálculo de módulo e fase.
        \end{itemize}
      \end{block}
    \end{column}
    \begin{column}{0.31\textwidth}
      \begin{block}{2. Simulação}
        \begin{itemize}
          \item circuito reproduzido no LTspice;
          \item alimentação simétrica de $\pm10\,\mathrm{V}$;
          \item senoide de $100\,\mathrm{mV}$;
          \item varredura de $40\,\mathrm{Hz}$ a $11\,\mathrm{kHz}$;
        \end{itemize}
      \end{block}
    \end{column}
    \begin{column}{0.31\textwidth}
      \begin{block}{3. Ensaio em bancada}
        \begin{itemize}
          \item montagem em protoboard;
          \item amplificador operacional e fonte simétrica;
          \item gerador e osciloscópio;
          \item conferência dos componentes;
          \item dois canais de medição;
        \end{itemize}
      \end{block}
    \end{column}
  \end{columns}

  \vspace{0.08cm}
  \begin{alertblock}{Processamento dos sinais}
    \centering
    $|H(j\omega)|=V_o/V_i$
    \qquad e \qquad
    $\phi=360^{\circ}\,\Delta T/T$.
  \end{alertblock}
\end{frame}